% ==============================================================================
% CONFIGURACIÓN IDIOMA, FUENTES Y TEXTO
% ==============================================================================
\usepackage[spanish,es-lcroman]{babel}   % Lowcase Roman: i), ii) en vez de I), II)
\renewcommand{\spanishtablename}{Tabla}

\usepackage[T1]{fontenc}       % Codificación de salida:  permite copiar/pegar acentos y eñes en el PDF
\usepackage[utf8]{inputenc}    % Codificación de entrada: permite escribir acentos y eñes directo desde el teclado

           
\usepackage{CJKutf8}           % Para japonés porque soy un friki

\usepackage{textcomp}          % Soporte para símbolos especiales en texto plano (euro, grados, etc.)

% ==============================================================================
% ENTORNOS MATEMÁTICOS Y SÍMBOLOS
% ==============================================================================
\usepackage{amsmath} % Para comandos de estructura matemática como \align y demás
\usepackage{amssymb} % Para símbolos matemáticos: \mathbb
\usepackage{bbm}     % Para más símbolos matemáticos: \mathbbm

\usepackage{amsthm}  % Para gestionar las definiciones y teoremas

% ==============================================================================
% DISEÑO DE PÁGINA, MARGENES Y ESTILOS
% ==============================================================================
\usepackage[dvipsnames, table]{xcolor}     % Para colores
\usepackage[most]{tcolorbox} % Añadir al principio del preámbulo
\usepackage{geometry}    % Dimensiones de página y margen de encuadernación (añadir 6mm si eso al terminar)
\geometry{top=25mm, bottom=25mm, left=15mm, right=15mm, bindingoffset=0mm}
\setlength{\headsep}{20pt}      % Distancia encabezadoi al texto

\usepackage{fancyhdr}           % Personalización avanzada de encabezados y pies de página
\setlength{\headheight}{15pt}   % El error del encabezado pequeño

% Para poner el interlineado que quiera
\usepackage{setspace}
\setstretch{1.1} 

% Para poner sangría
\setlength{\parindent}{0pt}

% Para modificar arrays y tablas
\usepackage{array}
\usepackage{longtable}
\usepackage{booktabs}

% Enumeración (como el itemize)
\usepackage{enumitem}

% Revisión, subrayado
\usepackage{soul}
\setulcolor{red}

% ==============================================================================
% Estilo de los teoremas y secciones
% ==============================================================================

% ==============================================================================
% Estilo de los teoremas y secciones
% ==============================================================================
\usepackage{titlesec}

% Ajusta el formato de \section: {formato}{etiqueta}{separación}{código antes}{código después}
% Estructura: \titlespacing*{comando}{margen izquierdo}{espacio antes}{espacio después}

% --- SECCIÓN: Centrada, grande y numerada ---
\titleformat{\section}[block]
{\normalfont\Large\bfseries\color{black}\centering}
{\thesection.}
{0.5em}
{}
\titlespacing*{\section}{0pt}{*2}{*3}

% --- SUBSECCIÓN: A la izquierda, normal y SIN numeración ---
\titleformat{s}
{\normalfont\normalsize\bfseries\color{black}\raggedright}
{}
{0em}
{}
\titlespacing*{\subsection}{0pt}{*2}{*1}

% --- SUBSUBSECCIÓN: Más grande (ej. \large) y a la izquierda ---
\titleformat{\subsubsection}
{\normalfont\large\bfseries\color{black}\raggedright} % Cambia \large por \Large si la quieres aún más grande
{\thesubsubsection} % Pon {} aquí si tampoco quieres que esté numerada
{0.5em}
{}
\titlespacing*{\subsubsection}{0pt}{*1.5}{*1}

% --- 2. CONFIGURACIÓN DE LAS SUBSECCIONES 

\newenvironment{enunciado}
  {\par\large\color{black}}
  {\par\medskip}

% Entorno para la resolución (Azul)
\newenvironment{solucion}
  {\par\large\color{blue}}
  {\par\bigskip}

\newtheoremstyle{Proposición}
  {5pt}         % Espacio arriba
  {10pt}        % Espacio abajo
  {\normalfont} % Fuente del cuerpo
  {}            % Sangría del cuerpo 
  {\bfseries}   % Fuente del encabezado (Definición 3.1)
  {.}           % Puntuación después del encabezado
  {.5em}           % Espacio después del encabezado
  {\thmname{#1}\thmnumber{ #2}\thmnote{ (#3)}} % Formato del encabezado

\newtheoremstyle{Ejemplo}
  {5pt}         % Espacio arriba
  {10pt}        % Espacio abajo
  {\normalfont} % Fuente del cuerpo
  {}            % Sangría del cuerpo 
  {\bfseries}   % Fuente del encabezado (Definición 3.1)
  {.}           % Puntuación después del encabezado
  {.5em}           % Espacio después del encabezado
  {Ejemplo \thmnumber{#2}} % Formato del encabezado

\newtheoremstyle{Ejercicio}
  {5pt}         % Espacio arriba
  {10pt}        % Espacio abajo
  {\normalfont} % Fuente del cuerpo
  {0pt}            % Sangría del cuerpo 
  {\bfseries}   % Fuente del encabezado 
  {}           % Puntuación después del encabezado
  {0pt}           % Espacio después del encabezado
  {Ejercicio \thmnumber{#2}\newline} % Formato del encabezado

\newtheoremstyle{Ampliacion}
  {5pt}         % Espacio arriba
  {10pt}        % Espacio abajo
  {\normalfont} % Fuente del cuerpo
  {0pt}            % Sangría del cuerpo 
  {\bfseries}   % Fuente del encabezado 
  {}           % Puntuación después del encabezado
  {0pt}           % Espacio después del encabezado
  {\centering Complemento \newline} % Formato del encabezado


%-----------------------------------------------------------------------------------------------
%-----------------------------------------------------------------------------------------------

\theoremstyle{Proposición}
\newtheorem*{prop}{Proposición}

\tcolorboxenvironment{prop}{
  enhanced,
  breakable,                % Permite que el ejemplo salte de página si es muy largo
  colback=black!5!white,    % Fondo verde muy claro (5% verde, 95% blanco)
  colframe=black!50!black,  % Borde verde oscuro
  boxrule=0.5pt,            % Grosor del borde
  arc=3pt,                  % Esquinas redondeadas
  top=2pt, bottom=2pt,      % Espaciado interior arriba y abajo
  left=5pt, right=5pt       % Espaciado interior a los lados
}

\newtheorem*{deff}{Definición}
\tcolorboxenvironment{deff}{
  enhanced,
  breakable,                % Permite que el ejemplo salte de página si es muy largo
  colback=black!5!white,    % Fondo verde muy claro (5% verde, 95% blanco)
  colframe=black!50!black,  % Borde verde oscuro
  boxrule=0.5pt,            % Grosor del borde
  arc=3pt,                  % Esquinas redondeadas
  top=2pt, bottom=2pt,      % Espaciado interior arriba y abajo
  left=5pt, right=5pt       % Espaciado interior a los lados
}

\theoremstyle{Ejemplo}
\newtheorem{ej}{Ejemplo}
\tcolorboxenvironment{ej}{
  enhanced,
  breakable,                % Permite que el ejemplo salte de página si es muy largo
  colback=green!5!white,    % Fondo verde muy claro (5% verde, 95% blanco)
  colframe=green!50!black,  % Borde verde oscuro
  boxrule=0.5pt,            % Grosor del borde
  arc=3pt,                  % Esquinas redondeadas
  top=2pt, bottom=2pt,      % Espaciado interior arriba y abajo
  left=5pt, right=5pt       % Espaciado interior a los lados
}

\theoremstyle{Ejercicio}
\newtheorem{ex}{}
\tcolorboxenvironment{ex}{
  enhanced,
  breakable,                % Permite que el ejemplo salte de página si es muy largo
  colback=red!5!white,    % Fondo verde muy claro (5% verde, 95% blanco)
  colframe=red!50!black,  % Borde verde oscuro
  boxrule=0.5pt,            % Grosor del borde
  arc=3pt,                  % Esquinas redondeadas
  top=2pt, bottom=2pt,      % Espaciado interior arriba y abajo
  left=5pt, right=5pt       % Espaciado interior a los lados
}


\theoremstyle{Ejemplo}
\newtheorem{_ejs}[ej]{Ejemplo}
\newenvironment{ejs}
  {\begin{_ejs}\leavevmode\vspace{-0.2\baselineskip}\begin{itemize}[leftmargin=15pt]}
  {\end{itemize}\end{_ejs}}

\tcolorboxenvironment{ejs}{
  enhanced,
  breakable,                % Permite que el ejemplo salte de página si es muy largo
  colback=green!5!white,    % Fondo verde muy claro (5% verde, 95% blanco)
  colframe=green!50!black,  % Borde verde oscuro
  boxrule=0.5pt,            % Grosor del borde
  arc=3pt,                  % Esquinas redondeadas
  top=2pt, bottom=2pt,      % Espaciado interior arriba y abajo
  left=5pt, right=5pt       % Espaciado interior a los lados
}


\newenvironment{amp}
  {\begin{center}
  \large\textbf{Ampliación}
  \end{center}}
  {}

\tcolorboxenvironment{amp}{
  enhanced,
  breakable,                % Permite que el ejemplo salte de página si es muy largo
  colback=blue!5!white,    % Fondo verde muy claro (5% verde, 95% blanco)
  colframe=blue!50!black,  % Borde verde oscuro
  boxrule=0.5pt,            % Grosor del borde
  arc=3pt,                  % Esquinas redondeadas
  top=2pt, bottom=2pt,      % Espaciado interior arriba y abajo
  left=5pt, right=5pt       % Espaciado interior a los lados
}

% ==============================================================================
% IMÁGENES, TABLAS Y LEYENDAS
% ==============================================================================
\usepackage{graphicx}         % Permite insertar imágenes (\includegraphics) y manipular su escala o rotación
\usepackage{float}            % Añade el modificador [H] para fijar la posición exacta de las imágenes
\usepackage{caption}          % Control personalizado de los títulos (leyendas) de figuras y tablas
\usepackage{subcaption}       % Permite subfiguras (ej. Figura 1a, 1b) dentro de un entorno figure
\graphicspath{ {imagenes/} }  % Para decirle a graphics que busque a en la carptea imagenes

% ==============================================================================
% TIKZ PARA PINTAR
% ==============================================================================
\usepackage{tikz}
\usetikzlibrary{calc}
\usetikzlibrary{arrows}
\usetikzlibrary{positioning}
\usetikzlibrary {arrows.meta}
\usetikzlibrary{decorations.markings}
\usetikzlibrary{hobby}
\usepackage{tikz-cd}
\usetikzlibrary{babel}
\usetikzlibrary{positioning, calc, fit, decorations.pathreplacing, arrows.meta}


% ==============================================================================
% MICROTYPE
% ==============================================================================
\usepackage{microtype} % Para mejorar estéticamente la composición (bastante)

% ==============================================================================
% GESTIÓN DE ENLACES Y METADATOS DEL PDF
% ==============================================================================
\usepackage[pdftex,     
                unicode,        % Para poder escribir en japonés           
                pdfauthor={     % Autor en los metadatos del archivo PDF
                        ジェクザ
                    },          
                pdftitle={      % Título en los metadatos del archivo PDF
                    Temario práctico de topología
                    }
            ]{hyperref} % Tema en los metadatos del archivo PDF

% ==============================================================================
% Comandos personalizados
% ==============================================================================

% ==============================================================================
% Comandos personalizados
% ==============================================================================

\newcommand{\R}{\mathbb{R}}
\newcommand{\Q}{\mathbb{Q}}
\newcommand{\Irr}{\mathbb{I}}
\newcommand{\C}{\mathbb{C}}
\newcommand{\K}{\mathbb{K}}
\newcommand{\N}{\mathbb{N}}
\newcommand{\Z}{\mathbb{Z}}
\newcommand{\F}{\mathbb{F}}
\newcommand{\cT}{\mathcal{T}}
\newcommand{\cS}{\mathcal{S}}
\newcommand{\cA}{\mathcal{A}}
\newcommand{\cB}{\mathcal{B}}
\newcommand{\cC}{\mathcal{C}}
\newcommand{\cD}{\mathcal{D}}
\newcommand{\cE}{\mathcal{E}}
\newcommand{\cF}{\mathcal{F}}
\newcommand{\cP}{\mathcal{P}}

\newcommand{\esptop}{(X,\mathcal{T})}
\newcommand{\topusu}{(\mathbb{R},\mathcal{T}_u)}

\newcommand{\seatop}{\text{Sea } (X,\mathcal{T}) \text{ un espacio topológico}}

\newcommand{\tn}[1]{\normalfont{#1}}
\newcommand{\tit}[1]{\textit{#1}}
\newcommand{\tbf}[1]{\textbf{#1}}
\newcommand{\titbf}[1]{\textit{\textbf{#1}}}
\newcommand{\tbfn}[1]{{\normalfont\textbf{#1}}}
\newcommand{\gr}[1]{\text{gr}(#1)}
\newcommand{\mcm}[1]{\text{mcm}(#1)}
\newcommand{\mcd}[1]{\text{mcd}(#1)}
\newcommand{\eqmod}[3]{#1 \equiv #2 \text{ (mod }#3)}
\newcommand{\neqmod}[3]{#1 \not\equiv #2 \text{ (mod }#3)}
\newcommand{\modulo}[2]{#1 \text{ mod } #2}
\newcommand{\im}[1]{\text{im}(#1)}
\newcommand{\evb}{\operatorname{ev}_\beta}
\newcommand{\gen}[1]{\langle#1\rangle}


\newcommand{\norm}[1]{\|#1\|}
\newcommand{\inv}[1]{#1^{-1}}
\newcommand{\abs}[1]{\left|#1\right|}

%Puntos de un subconjunto
\newcommand{\interior}[1]{\operatorname{int}(#1)}
\newcommand{\exterior}[1]{\operatorname{ext}(#1)}
\newcommand{\frontera}[1]{\operatorname{front}(#1)}


\newcommand{\wpar}[1]{\left(#1\right)}
\newcommand{\wbrac}[1]{\left\{#1\right\}}
\newcommand{\wcor}[1]{{\left[#1\right]}}

% Para recordarme mis comentarios
\newcommand{\warning}{%
  \PackageWarning{autor}{***** WARNING: REVISAR *****}%
}


 % Tamaño cabecera y pie de página
\renewcommand{\headrulewidth}{0.3pt}
\renewcommand{\footrulewidth}{0.3pt}


